\section*{Problemas}

\subsection*{Enunciados de los problemas}

% Recordar
\begin{problema}
	\label{prob:85eb17b}
	Un modelo cuadrático ajustado es $\hat{\texttt{mpg}} = 64.20 - 0.7634\,\texttt{hp} + 0.003657\,\texttt{hp}^2$. Calcule el \texttt{mpg} predicho para automóviles con $\texttt{hp}=60$, $\texttt{hp}=100$, y $\texttt{hp}=180$, e identifique en qué valor aproximado de \texttt{hp} el modelo predice el \texttt{mpg} mínimo.
\end{problema}

% Comprender
\begin{problema}
	\label{prob:5afaae8}
	Al ajustar dos modelos sobre el mismo conjunto de datos (\texttt{hp} vs. \texttt{mpg}), se obtiene $R^2_{\text{lineal}}=0.080$ para el modelo $\texttt{mpg}=c_0+c_1\,\texttt{hp}$ y $R^2_{\text{cuadrático}}=0.994$ para el modelo $\texttt{mpg}=c_0+c_1\,\texttt{hp}+c_2\,\texttt{hp}^2$. Interprete esta diferencia dramática y explique qué le dice sobre la verdadera forma funcional de la relación entre las variables.
\end{problema}

% Aplicar
\begin{problema}
	\label{prob:e757ebc}
	Con $n=5$ datos de \texttt{hp} y \texttt{mpg}, se obtuvo $R^2_{\text{lineal}}=0.080$ y $R^2_{\text{cuadrático}}=0.994$. Realice la Prueba $F$ Parcial ($H_0:c_2=0$, es decir, el término cuadrático no es necesario) al $\alpha=0.05$ (cuantil crítico $F_{0.05,1,2}=18.51$), usando $F_{\text{parcial}}=\dfrac{(R^2_{\text{cuad}}-R^2_{\text{lin}})/(2-1)}{(1-R^2_{\text{cuad}})/(5-2-1)}$.
\end{problema}

% Analizar
\begin{problema}
	\label{prob:c772cd1}
	Demuestre que el problema de ajustar por mínimos cuadrados un modelo polinomial $Y=c_0+c_1X+c_2X^2+\varepsilon$ es formalmente idéntico al problema de regresión lineal múltiple si se define la variable auxiliar $Z=X^2$ y se ajusta $Y=c_0+c_1X+c_2Z+\varepsilon$. Explique por qué la Ecuación Normal $\hat{\mathbf{c}}=(\mathbf{X}^T\mathbf{X})^{-1}\mathbf{X}^T\mathbf{Y}$ se aplica sin ninguna modificación, a pesar de que el modelo es \emph{no lineal} en la variable original $X$.
\end{problema}

% Evaluar
\begin{problema}
	\label{prob:6954b16}
	\textbf{La multicolinealidad inherente de los modelos polinomiales:} en un modelo polinomial $Y=c_0+c_1X+c_2X^2+\varepsilon$ donde $X$ toma valores estrictamente positivos concentrados en un rango angosto (por ejemplo, $X\in[95,105]$), demuestre que las columnas $X$ y $X^2$ de la matriz de diseño están necesariamente muy correlacionadas entre sí (correlación cercana a $1$), y evalúe por qué esto infla el VIF de ambos predictores sin que esto invalide las predicciones del modelo, aunque sí complica la interpretación aislada de $c_1$ y $c_2$.
\end{problema}

% Crear
\begin{problema}
	\label{prob:bc664de}
	Diseñe un problema original de regresión polinomial (distinto de la relación \texttt{hp}-\texttt{mpg} ya vista): un economista de precios modela las unidades vendidas ($Y$) de un producto en función de su precio de venta ($X$, en dólares), observando que la relación no es monótona: ventas relativamente altas a precios bajos, ventas mínimas a precios intermedios, y un repunte a precios muy altos (percepción de exclusividad). Se ajusta $\hat Y = 850 - 12X + 0.09X^2$. Calcule las ventas predichas para $X=30$, $X=70$ y $X=110$, e identifique el precio donde el modelo predice el mínimo de ventas.
\end{problema}

\subsection*{Soluciones detalladas}

% Recordar
\begin{solproblema}[prob:85eb17b]
	\begin{align}
		\hat{\texttt{mpg}}(60) &= 64.20-0.7634(60)+0.003657(60)^2 \approx 31.57.\\
		\hat{\texttt{mpg}}(100) &= 64.20-0.7634(100)+0.003657(100)^2 \approx 24.43.\\
		\hat{\texttt{mpg}}(180) &= 64.20-0.7634(180)+0.003657(180)^2 \approx 45.32.
	\end{align}
	El vértice de la parábola (mínimo) ocurre en $\texttt{hp}^*=-c_1/(2c_2)=0.7634/(2\times0.003657)\approx104.4$, es decir, el \texttt{mpg} mínimo se predice alrededor de $\texttt{hp}\approx104$.
\end{solproblema}

% Comprender
\begin{solproblema}[prob:5afaae8]
	La diferencia de $R^2=0.080$ a $R^2=0.994$ es dramática y revela que la relación entre \texttt{hp} y \texttt{mpg} \textbf{no es lineal en absoluto}: el modelo lineal apenas explica el $8\%$ de la variabilidad, prácticamente indistinguible de ruido, mientras que el modelo cuadrático explica el $99.4\%$. Esto confirma que la verdadera forma funcional subyacente es cuadrática (o al menos que un término de curvatura es indispensable), y que cualquier conclusión basada únicamente en el modelo lineal sería completamente engañosa.
\end{solproblema}

% Aplicar
\begin{solproblema}[prob:e757ebc]
	\begin{align}
		F_{\text{parcial}} = \frac{(0.994-0.080)/(2-1)}{(1-0.994)/(5-2-1)} = \frac{0.914/1}{0.006/2} = \frac{0.914}{0.003} \approx304.7.
	\end{align}
	Como $304.7\gg18.51$, \textbf{se rechaza $H_0$} de forma contundente: el término cuadrático $c_2\,\texttt{hp}^2$ es altamente significativo y necesario en el modelo; el modelo lineal simple está gravemente mal especificado para estos datos.
\end{solproblema}

% Analizar
\begin{solproblema}[prob:c772cd1]
	Definiendo $Z=X^2$ como una variable auxiliar calculada explícitamente, el modelo $Y=c_0+c_1X+c_2X^2+\varepsilon$ se reescribe exactamente como $Y=c_0+c_1X+c_2Z+\varepsilon$, que es \emph{lineal en los parámetros} $c_0,c_1,c_2$ —la única condición que requiere la teoría de mínimos cuadrados, independientemente de si las variables predictoras ($X$ y $Z=X^2$) tienen o no una relación no lineal entre sí. La matriz de diseño simplemente incluye una columna adicional con los valores de $X^2$ ya calculados: $\mathbf{X}_{\text{diseño}}=[\mathbf{1},\ X,\ X^2]$. Como la Ecuación Normal $(\mathbf{X}_{\text{diseño}}^T\mathbf{X}_{\text{diseño}})^{-1}\mathbf{X}_{\text{diseño}}^T\mathbf{Y}$ solo exige que la relación entre $Y$ y los parámetros $\mathbf{c}$ sea lineal (no que la relación entre $Y$ y $X$ lo sea), se aplica exactamente igual sin ninguna modificación.
\end{solproblema}

% Evaluar
\begin{solproblema}[prob:6954b16]
	Si $X$ toma valores concentrados en un rango angosto alrededor de un valor central grande (por ejemplo $X_0=100$), podemos escribir $X=X_0+\delta$ con $|\delta|$ pequeño relativo a $X_0$. Entonces:
	\begin{align}
		X^2 = X_0^2+2X_0\delta+\delta^2 \approx X_0^2+2X_0\delta \quad (\delta^2 \text{ despreciable frente a } 2X_0\delta \text{ cuando } |\delta|\ll X_0),
	\end{align}
	de modo que $X^2$ es aproximadamente una función \textbf{lineal} de $X$ (con pendiente $2X_0$) en ese rango angosto, lo que produce una correlación muestral $r_{X,X^2}$ muy cercana a $1$. Por la identidad $\text{VIF}_j=1/(1-R_j^2)$, una correlación tan alta entre $X$ y $X^2$ produce un $R_j^2$ cercano a $1$ al regredir una contra la otra, y por tanto un VIF muy elevado para ambos predictores.

	Sin embargo, esto \textbf{no invalida las predicciones} $\hat Y$ del modelo, que dependen de la combinación conjunta $\hat c_1X+\hat c_2X^2$ y permanecen numéricamente estables y precisas. Lo que sí se ve comprometido es la interpretación marginal aislada de $\hat c_1$ o $\hat c_2$ individualmente (sus errores estándar se inflan, y separar el efecto de $X$ del efecto de $X^2$ pierde sentido práctico), exactamente el mismo fenómeno interpretativo discutido para predictores correlacionados en general.
\end{solproblema}

% Crear
\begin{solproblema}[prob:bc664de]
	\begin{align}
		\hat Y(30) &= 850-12(30)+0.09(30)^2 = 850-360+81 = 571 \text{ unidades.} \\
		\hat Y(70) &= 850-12(70)+0.09(70)^2 = 850-840+441 = 451 \text{ unidades.} \\
		\hat Y(110) &= 850-12(110)+0.09(110)^2 = 850-1320+1089 = 619 \text{ unidades.}
	\end{align}
	El vértice de la parábola (mínimo de ventas) ocurre en:
	\begin{align}
		X^* = -\frac{c_1}{2c_2} = \frac{12}{2(0.09)} = \frac{12}{0.18} \approx 66.67,
	\end{align}
	es decir, el modelo predice que las ventas alcanzan su punto más bajo alrededor de un precio de $\$66.67$, consistente con el valor mínimo observado entre los tres precios evaluados ($\hat Y(70)=451$, el más cercano al vértice).
\end{solproblema}
